\section{Conclusion}

\label{sec:conclusion}

This study makes authorization scope an explicit experimental factor in runtime guardrail evaluation. Corrected target matching separates security-target scope from complete action arguments, while paired comparisons connect authorization information with benign completion and operational burden.

The 900-assignment study shows improved AgentDojo completion with exact-target assistance. The 480-assignment contemporaneous calibration identifies a marked utility cost under compiled authorization, with an uncertain security contrast. The trusted-intent review locates compiler-coverage and clarification needs, and sixteen synthetic controls distinguish policy from full-module flow behavior within the tested scenarios.

The empirical value is a clearer account of the tested security--utility trade-off, rather than a general proof of full-module superiority. Keeping authorization, Guard decisions, execution, completion, failures, and unknowns separate provides a practical basis for subsequent non-oracle runtime evaluations.
